\chapter{User Documentation}
\label{app:user}

This chapter describes how to run the application and how to provide input data.

\section{Running the Application}

The application can be executed in two ways:
\begin{itemize}
    \item using a graphical user interface (GUI),
    \item using a configuration file.
\end{itemize}

\subsection{Running with GUI}

If the application is started without command-line arguments, a graphical interface is launched.

The procedure is as follows:
\begin{enumerate}
    \item Run the application.
    \item Select the input JSON file.
    \item Configure packing settings.
    \item Select the evolutionary algorithm.
    \item Set the algorithm hyperparameters.
    \item Start the computation.
\end{enumerate}

\subsection{Running with Configuration File}

The application can also be executed from the command line using a configuration file:

\begin{verbatim}
dotnet run --project src/App config.json
\end{verbatim}

The file \texttt{config.json} contains all required settings.

\section{Configuration Parameters}

The program is configured using a \texttt{ProgramSetting} structure, which represents the full high-level configuration of the application.

Conceptually, the configuration has the following structure:

\begin{verbatim}
ProgramSetting
{
    SourceFile: string,
    OutputFile: string,
    EvolutionStatisticsFile: string?,

    PackingSetting:
    {
        SelectedPlacementHeuristics: string[],
        AllowRotations: bool,
        PackingOrder: string? or null
    },

    EvolutionAlgorithmSetting:
    {
        (common parameters)
        TargetFitness: double?,
        Individuals: int,
        UseIndividualsAsRelative: bool,
        NumberOfGenerations: int,

        (algorithm-specific parameters)
        ...
    }
}
\end{verbatim}

\subsection{Input and Output}

These parameters correspond to the top-level fields of \texttt{ProgramSetting}:

\begin{itemize}
    \item \texttt{SourceFile} -- path to the input JSON file containing the packing problem,
    
    \item \texttt{OutputFile} -- path where the resulting packing solution will be saved,
    
    \item \texttt{EvolutionStatisticsFile} -- optional CSV file where evolution statistics are exported.
\end{itemize}

\subsection{Packing Settings}

These parameters belong to the \texttt{PackingSetting} structure and define the packing constraints.

\begin{itemize}
    \item \texttt{SelectedPlacementHeuristics} -- list of heuristics used to determine where a box should be placed,
    
    \item \texttt{AllowRotations} -- allows boxes to be rotated when searching for a valid placement,
    
    \item \texttt{PackingOrder} -- optional heuristic defining the order in which boxes are packed.
\end{itemize}

If \texttt{PackingOrder} is not specified, the order of boxes is determined by the evolutionary algorithm.

\subsection{Evolutionary Algorithm Settings}

These parameters belong to the \texttt{EvolutionAlgorithmSetting} structure.
The available parameters depend on the selected algorithm.

\begin{itemize}
    \item \texttt{AlgorithmName} -- discriminator that determines which evolutionary algorithm configuration is used,

    \item \texttt{TargetFitness} -- optional target fitness value for early stopping,
    
    \item \texttt{Individuals} -- number of individuals in the population, denoted as $M$ in the algorithm,

    \item \texttt{UseIndividualsAsRelative} -- indicates whether \texttt{Individuals} is interpreted as a relative value; if \texttt{true}, the actual population size is scaled proportionally to the input size,

    \item \texttt{NumberOfGenerations} -- maximum number of generations, corresponding to $G$.
\end{itemize}

\subsubsection{Elitist Genetic Algorithm Settings}

\begin{itemize}
    \item \texttt{PercentageOfEliteIndividuals} -- percentage of elite individuals preserved in each generation, corresponding to $p_E$,
    
    \item \texttt{PercentageOfElementsFromElite} -- percentage of elements inherited from the elite parent during crossover, corresponding to $p_C$,
    
    \item \texttt{PercentageOfElementsMutated} -- percentage of elements mutated in each offspring, corresponding to $p_M$,
    
    \item \texttt{PercentageOfRandomIndividuals} -- percentage of new random individuals introduced each generation, corresponding to $p_R$,
    
    \item \texttt{HillClimbingIterations} -- number of hill climbing iterations applied to each offspring, corresponding to $I_{HC}$,
    
    \item \texttt{HillClimbingPercentageOfElementsMutated} -- mutation rate used during hill climbing, corresponding to $p_{HC}$.
\end{itemize}

\subsubsection{Probabilistic Hill Climbing Settings}

\begin{itemize}
    \item \texttt{PercentageOfElementsChanged} -- percentage of elements modified when generating a neighboring solution, corresponding to $p_M$,
    
    \item \texttt{StartAcceptanceProbability} -- probability of accepting a worse solution at the beginning of the search, corresponding to $p_A^{(0)}$,
    
    \item \texttt{EndAcceptanceProbability} -- minimum probability of accepting a worse solution, corresponding to $p_{\min}$,
    
    \item \texttt{AcceptanceDecayFactor} -- factor by which the acceptance probability decreases each generation, corresponding to $\alpha$.
\end{itemize}

\subsection{Configuration Schema and Examples}

The JSON schemas defining the configuration structure are located in the \texttt{schema} directory:

\begin{itemize}
    \item \texttt{schema/ProgramSetting.json} -- schema for the high-level application configuration,
    \item \texttt{schema/PackingInput.json} -- schema for the packing problem input data.
\end{itemize}

Example input and output files are available in the \texttt{examples} directory.

\section{Packing Input Format}

The packing input is provided as a JSON file describing container properties and a list of boxes.

\begin{verbatim}
{
  "ContainerProperties": {
    "Sizes": { "X": 100, "Y": 100, "Z": 100 },
    "MaxWeight": 500000
  },
  "BoxPropertiesList": [
    { "ID": 0, "Sizes": { "X": 20, "Y": 10, "Z": 5 }, "Weight": 123 }
  ]
}
\end{verbatim}

\section{Output}

The program produces:
\begin{itemize}
    \item a JSON file with the resulting packing solution,
    \item optionally a CSV file with evolution statistics.
\end{itemize}